\documentclass[../../syllabus_2.tex]{subfiles}
\begin{document}
\chapter{Proportionnalité, projections parallèles}

\section{Tableau de proportionnalité}

\begin{exercise}
    Pour son anniversaire, Diego reçoit un GSM. Il lui reste à choisir un abonnement. Il se rend chez un opérateur qui lui propose deux tarifs : le tarif standard à 0,50~€/minute d'appel et SMS illimités ou le tarif premium à 10~€ fixe et 0,25~€/minute d'appel et SMS illimités.
    \begin{enumerate}
        \item Pour chaque tarif, établis un tableau : \\
              Tarif standard :
              \begin{center}
                  \begin{tblr}{
                          hlines,
                          vlines,
                          columns={c},
                          rows={m},
                          column{1}={5cm,l}
                      }
                      Minutes appelées  & 0 & 5 & 10 & 15 & 20 & 25 & 30 & 35 \\
                      Prix à payer en € &   &   &    &    &    &    &    &
                  \end{tblr}
              \end{center}

              Tarif premium:
              \begin{center}
                  \begin{tblr}{
                          hlines,
                          vlines,
                          columns={c},
                          rows={m},
                          column{1}={5cm,l}
                      }
                      Minutes appelées  & 0 & 5 & 10 & 15 & 20 & 25 & 30 & 35 \\
                      Prix à payer en € &   &   &    &    &    &    &    &
                  \end{tblr}
              \end{center}

        \item Représente graphiquement ces deux tarifs ci-dessous.
              \begin{center}
                  \includegraphics[width=\linewidth]{chapitres/prop-et-proj/tikz/tikz001.tikz}
              \end{center}

        \item Ces deux tarifs représentent-ils des relations directement proportionnelles ? Justifie.
              \adf{2}

              \hskip -20pt
              \adf[\textbf{Définition :} Deux grandeurs sont proportionnelles si ]{3}
              % on obtient la seconde grandeur en multipliant la première par un nombre constant. Ce nombre est appelé le coefficient de proportionnalité et est noté k.


        \item Si on pose \(x\) le nombre de minutes et y le prix en €, écrit la relation qui unit ces deux valeurs dans la situation de proportionnalité.
              \adf{1}

        \item Dans cet exemple, le coefficient de proportionnalité vaut :

              \vspace{1cm}

        %\item Comment peut-on reconnaitre graphiquement qu'il s'agit d'une relation directement proportionnelle ?
        %      \adf{2} --> question posée dans l'exercice 3
    \end{enumerate}
\end{exercise}
         
    Formule générale
    \begin{center}
    $ y = k \cdot x$
    \end{center}

    \begin{center}
    k = coefficient de proportionnalité = \(\dfrac{y}{x}= \dfrac{\text{seconde grandeur}}{\text{permière grandeur}}\)
    \end{center}

    Ainsi, lorsque nous nous retrouvons face à un tableau, nous pouvons déterminer le coefficient de proportionnalité en divisant la valeur dépendante par la valeur indépendante.
    \vspace{0.5cm}

        
    

%\newpage
\begin{exercise}
    Complète les tableaux suivants pour qu'ils correspondent à des situations de proportionnalité directe et indique le coefficient de proportionnalité \(k\) pour passer de \(x\) à  \(y\).
    \begin{enumerate}
        \item  \begin{tblr}[t]{hlines, vlines, columns={c, 1cm}, rows={m, 1cm}}
                  \(x\) & 5 & 8  & 10 & 16 \\
                  \(y\) &   & 48 & 66 &
              \end{tblr}
              \hspace{1cm}
              \(k\) =
        \item \begin{tblr}[t]{hlines, vlines, columns={c,1cm}, rows={m,1cm}}
                  \(x\) & 5 & 8  & 10 & 16 \\
                  \(y\) &   & 48 & 66 &
              \end{tblr}
              \hspace{1cm}
              \(k\) =
        \item \begin{tblr}[t]{hlines, vlines, columns={c,1cm}, rows={m,1cm}}
                  \(x\) & 5 & 8  & 10 & 16 \\
                  \(y\) &   & 48 & 66 &
              \end{tblr}
              \hspace{1cm}
              \(k\) =
    \end{enumerate}
\end{exercise}

\begin{exercise}
    Cristina et Elena, deux amies, vont à l'école à vélo. Elles partent au même moment de leur immeuble, mais ne font pas le trajet ensemble. Voici les données pour un matin du mois de mai.

    Pour Cristina :
    \begin{center}
        \begin{tblr}{hlines, vlines, columns={c}, rows={m, 1cm},column{1}={5cm,l}}
            Temps de trajet (minutes) & 0 & 5    & 10  & 15   & 20 & 25   & 30  \\
            Distance parcourue (km)   & 0 & 1,25 & 2,5 & 3,75 & 5  & 6,25 & 7,5
        \end{tblr}
    \end{center}

    Elena, elle, est arrivée en retard. Voici ce qu'elle explique à son éducateur : \textit{\og Je sais qu'il est 08 h 45 Monsieur, mais je n'ai pas eu de chance ce matin. Je suis partie de chez moi à 7 h 50, au début, j'allais à mon aise (10 km/h). A 8 h 05, je me suis arrêtée pour attendre Batuhan. Il est arrivé à 8 h 10, et comme on était déjà un peu en retard, on a accéléré (15km/h). Malheureusement, à 08 h 20 j'ai roulé sur un clou et j'ai crevé mon pneu. On a essayé de le réparer pendant 10 minutes, mais ça n'a pas marché. Je suis donc montée sur le porte-bagage de Batuhan, et il a pédalé pour deux. C'était plus dur (10 km/h). On est désolé Monsieur, on a fait de notre mieux.\fg}

    \begin{enumerate}
        \item Complète le tableau de valeurs pour le trajet d'Elena.
              \begin{center}
                  \begin{tblr}{hlines, vlines, columns={c,1cm}, rows={m, 1cm},column{1}={5cm,l}}
                      Temps de trajet (minutes) & 0 & 5 & 10 & 15 & 20 & 25 & 30 \\
                      Distance parcourue (km)   &   &   &    &    &    &    &
                  \end{tblr}
              \end{center}

              \newpage
        \item Trace le graphique qui représente le trajet d'Elena ainsi que celui de Cristina. Utilise des couleurs différentes.
              \begin{center}
                  \includegraphics[width=\linewidth]{chapitres/prop-et-proj/tikz/tikz002.tikz}
              \end{center}
        \item Quel graphique représente une relation proportionnelle entre distance et temps ?
              \adf{2}
        \item Quel est le coefficient de proportionnalité pour cette relation ?
              \adf{2}
    \end{enumerate}
\end{exercise}

\vspace{1cm}
\adf[Le graphique d'une relation proportionnelle entre deux grandeurs est toujours] {1} %une droite passant par l'origine.

\newpage
\section{Le produit en croix}

\adf[\textbf{Définition:} Une proportion est ]{1} %une égalité entre deux rapports.

\exemple On roule à la même vitesse si on parcours 120km en 3h ou 240km en 6h. 
Quelle est la vitesse en kilomètres-heure ? 
\vspace{1cm}
\\ 
\(\dfrac{120}{3}=\dfrac{240}{6}=\dfrac{\ldots\ldots}{\ldots\ldots}=\ldots\ldots km/h\)
\vspace{1cm}

Comme vu dans le chapitre sur les équations, lorsque nous avons une égalité de deux fractions  $\dfrac{a}{b}=\dfrac{c}{d}$, nous pouvons réécrire l'égalité comme étant:

\vspace{3cm}

%$a\cdot d = b \cdot c$. \\


\textbf{Vocabulaire:}Dans la proportion $\dfrac{a}{b}=\dfrac{c}{d}$,
\begin{enumerate}[label= ]
\item a, b, c et d sont les termes,
\item a et d sont les extrêmes et
\item b et c sont les moyens.
\end{enumerate}

\adf[\textbf{Propriété fondamentale:} Dans toute proportion,]{1} %le produit des extrêmes est égal au produit des moyens.

On peut donc échanger les moyens entre eux, et les extrêmes entre eux.

\begin{exercise}
    Si pour ton cocktail d'anniversaire, tu utilises dix parts de jus d'ananas et quinze parts de jus d'orange, quelle doit être la part de jus d'orange pour un cocktail aux mêmes proportions si tu n'as que huits parts de jus d'ananas?
\end{exercise}
\vspace*{5cm}

\begin{exercise}
    À un championnat de frisbee, on compte les points des équipes en fonction du nombre de buts marqués par rapport au nombre de tentatives. Deux équipes ont reçu le même nombre de points, la première équipe a marqué 12 goals, la deuxième équipe 10 goals sur 35 tentatives. Combien de tentatives la première équipe a-t-elle eues?  
\end{exercise}
\vspace*{5cm}

\begin{exercise}
    Nathan et Samuel veulent repeindre le local de leur classe en vert clair. Pour obtenir cette couleur, il faut mélanger du bleu et de jaune. Nathan a estimé qu'il fallait une portion de bleu pour cinq de jaune. Il effectue le mélange. À ce moment-là, Samuel arrive et lui dit "c'est beaucoup trop clair, il fallait faire le contraire : cinq portions de bleu et une de jaune". Combien Nathan doit-il rajouter de portions de bleu pour obtenir la bonne couleur ?
\end{exercise}
\vspace*{5cm}

\newpage
\section{Projections parallèles}
\subsection{Découverte}
Sara va se promener sur la plage avec son petit frère Juan. Sara mesure \SI{1,70}{m} et Juan \SI{1,20}{m}. Elle remarque que l'ombre de Juan mesure \SI{1,80}{m}. Quelle est la longueur de son ombre ?

\begin{center}
    \includegraphics[width=15cm]{chapitres/prop-et-proj/img/sun.jpg}
\end{center}



\vfill
Remarque :  Le fait que la projection est parallèle conserve les proportions entre les différentes mesures.

\newpage

\begin{exercise}
    Tu es parti en voyage à Tokyo. À 14 h 05, heure locale, tu remarques que ton ombre mesure \SI{85}{cm}. Tu mesures \SI{1,50}{m}. À côté de toi, s'étend l'ombre du Tokyo Skytree, un de plus grand gratte-ciel au monde. L'ombre mesure \SI{338}{m}. Ton frère te dit que ce gratte-ciel est plus grand que le Burj Khalifa à Dubaï qui mesure \SI{828}{m}. A-t-il raison ? \\ Argumente.
    %Faux, elle mesure 634m (338*8/15)
\end{exercise}
\vspace*{7cm}

\subsection{Définition}
Une projection parallèle est une transformation du plan qui projette un élément sur une droite de façon parallèle.
Pour construire l'image du point \(A\) par la projection sur la droite \(b\), parallèlement à la droite \(a\) :
\begin{itemize}
    \item Tracer la droite passant par \(A \) parallèle à \(a\).
    \item Pointer \(A'\) à l'intersection de cette parallèle et de \(b\).
          \begin{center}
              \includegraphics[width=11cm]{chapitres/prop-et-proj/img/projection.jpg}
          \end{center}
\end{itemize}

\begin{exercise}
    Projette le triangle \(ABC\), sur la droite \(d\), parallèlement à la droite \(a\). Laisse tes constructions visibles.
    \begin{center}
        \includegraphics[width=11cm]{chapitres/prop-et-proj/img/triangle.jpg}
    \end{center}
\end{exercise}

\begin{exercise}
    Te voici dans une situation délicate. Tu désires partager le segment \([AB]\) en 4 segments de même longueur sans utiliser le moindre instrument de mesure. Comment vas-tu faire ? Aide-toi des projections parallèles. \\

    \includegraphics[width=10cm]{chapitres/prop-et-proj/img/segment.jpg}
\end{exercise}

\newpage
Pour partager un segment en \(4\) parties égales, voici les étapes à suivre.

\adf[Etape 1: ]{2} %On trace une demi-droite d'origine A
\adf[Etape 2: ]{3} % Sur cette demi-droite on reporte 4 fois la même mesure à partir de A. On appelle les point obtenus respectivement C, D, E et F.
\adf[Etape 3: ]{1} %On relie F à B
\adf[Etape 4: ]{2} % On trace les parallèle à BF passant par C, D et E à l'aide de l'équerre

\begin{exercise}
    Partage le segment suivant en 6 segments de même longueur, sans rien mesurer. \\

    \includegraphics[width=10cm]{chapitres/prop-et-proj/img/segment.jpg}
\end{exercise}

\newpage
\section{Agrandissements et réductions}

Voici une image d'un chien et celle d'un ours. Complète la version réduite et agrandie de chacun, en respectant les proportions.
\begin{center}
    
\includegraphics[width=15cm]{chapitres/prop-et-proj/img/dog.jpg}
    
\includegraphics[width=15cm]{chapitres/prop-et-proj/img/bear.jpg}
\end{center}


Qu'as-tu besoin de calculer pour pouvoir respecter les proportions ?
\adf{1}
%le coefficient de prop

\newpage
\begin{exercise}
    Construis, dans le quadrillage suivant.
    \begin{enumerate}
        \item Réalise une réduction, $k=\dfrac{1}{2}$
              \medbreak
              \includegraphics[width=\linewidth]{chapitres/prop-et-proj/tikz/tikz003.tikz}

        \item Réalise un agrandissement, $k =3$
              \medbreak
              \includegraphics[width=\linewidth]{chapitres/prop-et-proj/tikz/tikz004.tikz}
    \end{enumerate}
\end{exercise}

\newpage
\textbf{Agrandissements et réductions}

Les agrandissements et les réductions conservent :
\begin{itemize}[label=$\bullet$]
    \item La forme des figures
    \item L'amplitude des angles
    \item Le parallélisme et la perpendicularité
\end{itemize}

Dans un agrandissement ou une réduction, les côtés homologues ont des longueurs proportionnelles :
\begin{itemize}[label=$\bullet$]
    \item Pour un agrandissement: \(k > 1\)
    \item Pour une réduction: \(k < 1\)
    \item Pour une isométrie : \(k = 1\)
\end{itemize}

\textbf{Calcul de périmètres et d'aires dans les homothéties}

La première figure est une réduction de facteur $\dfrac{1}{2}$ de la seconde figure. Calcule le périmètre et l'aire des deux figures.
\begin{center}
    \includegraphics[width=10cm]{chapitres/prop-et-proj/img/trapeze.jpg}
\end{center}

Périmètres :
\vspace{2cm}

Aires :
\vspace{2cm}

Que peux-tu conclure ? 
\adf{4}
%Pour deux figures agrandies (ou réduites) avec un coefficient k:
%- les longueurs (les périmètres) sont mulitpliés par k,
%- les aires sont multipliées par k².

Nous avons vu trois méthodes de résolution de problèmes de proportionnalité. Nous pouvons généralement choisir la méthode qui nous semble la plus adaptée à la situation. 


\exemple Marie a acheté 4 Mangas de OnePiece et elle a payé \SI{20}{euros}. Combien aurait-elle payé pour 6 albums ? Écris tous tes calculs.
\begin{enumerate}[itemsep=5cm]
    \item La règle de trois
    \item Le tableau de proportionnalité
    \item Le produit en croix
\end{enumerate}


\newpage
\section{Exercices mélangés}
\setcounter{exercise}{0}


\begin{exercise}
    Ma voiture consomme, en moyenne, 4 litres de carburant pour \SI{100}{km}. Quelle distance pourrais-je parcourir avec un plein de 50 litres ?
\end{exercise}

\begin{exercise}
    Chez Colruyt, un lot de 4 bouteilles de vin est affiché à \SI{22}{euros}. On peut acheter les bouteilles à la pièce. \\
    Combien vais-je payer pour l'achat de 10 bouteilles ?
\end{exercise}

\begin{exercise}
    Pour une exposition, un groupe de 12 élèves a payé \SI{84}{euros}. Qu'auraient payé 17 élèves ?
\end{exercise}

\begin{exercise}
    Sur un plan à l'échelle \(1/\num{2000}\), un champ rectangulaire est représenté par un rectangle mesurant \SI{2}{cm} de long et \SI{1,4}{cm} de large.
    \begin{enumerate}
        \item Quelles sont les dimensions réelles (en mètres)  de ce champ ?
        \item Quelle est l'aire de ce champ (en mètres carrés) ?
    \end{enumerate}
\end{exercise}

\begin{exercise}
    Pour télécharger 3 chansons sur internet, à une certaine époque, il fallait 1 minute.
    \begin{enumerate}
        \item  Complète, en te basant sur ce temps moyen de téléchargement, le tableau de proportionnalité suivant : (\emph{Sur cette feuille})\\
              \begin{center}
                  \begin{tblr}{hlines, vlines, columns={c}, rows={m, 1cm}}
                      Nombre de chansons & Durée du téléchargement \\
                                         & 120                     \\
                      9                  &                         \\
                                         & 500
                  \end{tblr}
              \end{center}

        \item  Calcule le nombre de chansons que tu pourrais télécharger, à la même vitesse, en une demi-heure.
    \end{enumerate}
\end{exercise}

\begin{exercise}
    Trace un semgent $[AT]$ de longueur quelconque. Ensuite, grâce à ce que tu as appris sur les projetions parallèles, place le point $S$ tel que $|AS|=\dfrac{3}{4}\cdot|AT|$
\end{exercise}
\end{document}